Recent anecdotal reports describe Large Language Models (LLMs) suggesting that users ``go to sleep'' or ``get some rest'' during prolonged interactions, particularly at night. It remains unclear whether this behavior is a universal artifact of interaction length (contextual decay) or a conditioned response to user state. We investigate this phenomenon using a multi-agent simulation framework, varying user personas (Neutral, Tired/Anxious, Energetic) and temporal contexts (2:00 PM vs. 3:00 AM). Our results show that the \studyname is overwhelmingly persona-conditioned: LLMs suggested sleep in 90\% of trials where the user exhibited a tired/anxious persona, but in 0\% of trials with neutral or energetic personas, regardless of the explicit time of day. While late-night context alone did not trigger the sleep suggestion, it acted as a significant accelerator when combined with a tired persona, reducing the average turn of suggestion from 5.22 to 3.89. These findings imply that the sleep suggestion is not a generic contextual failure but a targeted, empathetic adaptation to the perceived user state. This research has implications for AI alignment, showing that models can project affective states onto users based on subtle linguistic cues.
